%% kapa: exam=1 round=31 year=2020 period=1 form=A subject=civil_law subject_ko=민법 seq=3 number=3
%% source: Q-Net 감정평가사 기출문제 (cst00309) artlSeq=5207219 file=extracted/1008/5207219/A형/A형/2020년도 제31회 감정평가사 1차 시험_1교시_A형.pdf
%% conversion: kordoc 4.24.51 -> tools/gen1.py (auto); review=proofread-sonnet; flags=none
\begin{question}{3}
미성년자의 행위능력에 관한 설명으로 옳은 것은? (다툼이 있으면 판례에 따름)
\begin{choices}
  \item 행위능력제도는 자기책임의 원칙을 구현하여 거래의 안전을 도모하기 위한 것이다.
  \item 미성년자가 그 소유의 부동산을 그의 친권자에게 증여하고 소유권이전등기를 마친 경우, 다른 사정이 없으면 적법한 절차를 거친 등기로 추정된다.
  \item 친권자는 그의 미성년 자(子)의 이름으로 체결한 계약을 자(子)가 미성년임을 이유로 취소할 수 있다.
  \item 친권자가 그의 친구의 제3자에 대한 채무를 담보하기 위하여 미성년 자(子) 소유의 부동산에 담보를 설정하는 행위는 이해상반행위이다.
  \item 미성년자가 타인을 대리할 때에는 법정대리인의 동의를 얻어야 한다.
\end{choices}
\end{question}
\answer{2}
